\section{AI-agent usage in this analysis program}
\label{sec:ai_agent}

This internal note includes an explicit evaluation of AI-agent support for an
experimental QCD analysis workflow. The central question is not only whether an
agent can generate text or isolated scripts, but whether it can execute a full
analysis loop under physicist supervision: data and MC processing, unfolding
orchestration, systematic propagation, figure production, and note maintenance
with traceable links to code and outputs.

\subsection{Scope and operating model}

The agent used in this analysis is interfaced through Anthropic Claude Code and OpenAI Codex, depending on the preference of the physicist.
The operating model is human-in-the-loop: the physicist formulates goals and
provides expert direction through iterative prompts; the agent executes the
corresponding implementation steps; the physicist reviews the outputs and
provides feedback; the process is then repeated. Physics decisions remain
analyst-owned, while the agent provides implementation speed, consistency
checks, and documentation rigor. In practical terms, the agent is used to:
\begin{itemize}
\item execute and monitor multi-step workflows;
\item enforce consistent plotting and note formatting conventions;
\item connect analysis statements in the note to concrete script paths,
      configuration keys, and produced artifacts;
\item iterate quickly on uncertainty-view and validation presentation.
\end{itemize}

This model is intentionally conservative: the agent is treated as an
accelerator for engineering and documentation tasks, not as an autonomous scientific decision-maker at this stage. The interaction is analogous to advising a graduate student, where the physicist supplies expert judgment and direction while the student (agent) carries out the technical execution.

\subsection{Boundaries, controls, and reproducibility}

To keep scientific accountability explicit, each high-impact step has a
control point where analyst sign-off is required before results are treated
as baseline:
\begin{itemize}
\item event and object selection definition;
\item nominal unfolding regularization choice;
\item systematic-combination policy;
\item final figure and uncertainty-presentation choices;
\item interpretation language in the note.
\end{itemize}

Reproducibility is ensured through multiple complementary mechanisms. All
analysis steps are version-controlled, with code and configuration managed
in a git repository so that any result can be traced to a specific commit
and set of inputs. Workflow runs produce structured output artifacts,
including summary files that capture the configuration used, the software
versions, and checksums of key outputs. Physicist oversight at each review
stage provides an additional layer of validation: discrepancies between
expected and observed outputs are flagged before any result is promoted to
the baseline. These controls together ensure that results are reproducible
independently of the AI agent's involvement and that the physicist retains
full accountability for all scientific claims.

\subsection{Why this matters for the broader physics program}

The broader objective of this work is to extend modern AI, including large
language models (LLMs), beyond purely theoretical calculations in
fundamental physics. To fully close the theory--experiment loop, the
experimental side must produce reliable, uncertainty-aware observables and
covariance products that can be directly consumed by precision QCD fits.
In this vision, AI is applied not only to theoretical calculations but also
to experimental analysis workflows.

The goal is to accelerate progress in fundamental physics by leveraging AI
agents, LLMs, and other advanced models across both domains, ideally within
a feedback loop between theory and experiment. In this loop, illustrated in
Fig.~\ref{fig:ai_loop}, theory predictions inform experimental analyses,
whose results in turn constrain and guide further theoretical proposals.
Throughout this process, the physicist provides expert oversight and
validation at each stage.

Recent case studies in AI-assisted scientific discovery and HEP theory
support this direction, while also emphasizing verification,
reproducibility, and expert oversight as essential constraints
for reliable deployment of such systems
\cite{Woodruff:2026practicalsi,Bubeck:2025coscientist,OpenAITheory:2026singleminus,Schwartz:2026cparameterllm}.
In the present work, the experimental--theory bridge is represented by
release-grade covariance and correlation outputs for thrust, with the next
step being an expansion to additional event-shape observables such as the
C-parameter and heavy-jet mass, and ultimately to multidimensional
constraints.

This section therefore serves both as a methodological statement and as a proof of concept. In particular, it shows that an AI agent can
assist with end-to-end execution of experimental workflows when constrained
by explicit controls, validation checks, and human review. The operational
details of this collaboration model, including the interaction protocol and
persistent bottlenecks observed from the agent perspective, are documented
in Sec.~\ref{sec:challenges}.